%BeginMC
\question
Several species of warblers (a group of small colorful birds) in North America can live in the same tree only because they

\begin{choices}
	\choice created different types of nests within the tree.%EndChoice
	\choice eat different foods within the tree.%EndChoice
	\correctchoice partitioned resources to eliminate competition.%EndChoice
	\choice can fly to different heights within the tree.%EndChoice
	\choice out-competed other species of birds, such as sparrows.%EndChoice
\end{choices}
%EndMC


%BeginMC
\question
Competitive exclusion means that

\begin{choices}
	\correctchoice one species will always out compete another species that occupies the same niche.  %EndChoice
	\choice two species will partition resources to eliminate competition. %EndChoice
	\choice one individual will always out compete another individual in a population. %EndChoice
	\choice one population will always exclude another another population of the same species in a community. %EndChoice
	\choice a species will have a smaller realized niche than a fundamental niche. %EndChoice
\end{choices}
%EndMC

%BeginMC
\question
Cowbirds birds lay their eggs in the nests of other birds like cardinals. The cowbird chicks kick the cardinal chicks out of the nest, so the cardinal chicks die. The adult cardinals raise the cowbird chicks as their own. This is most likely an example of 

\begin{choices}
	\choice competition. %EndChoice
	\choice resource partitioning. %EndChoice
	\correctchoice parasitism. %EndChoice
	\choice symbiosis. %EndChoice
	\choice mutualism. %EndChoice
\end{choices}
%EndMC

%BeginMC
\question
Color that helps hide an animal in its environment is called

\begin{choices}
	\choice mimicry. %EndChoice
	\choice aposomatic coloration %EndChoice
	\correctchoice cryptic coloration. %EndChoice
	\choice warning coloration. %EndChoice
	\choice environmental coloration. %EndChoice
	\choice hideandseek coloration. %EndChoice
\end{choices}
%EndMC

%BeginMC
\question
Which block in the figure below is most consistent intermediate levels of disturbance in an ecosystem? (Different shades, including white, represent different species. The small squares indicate relative abundance of individuals present.)

\begin{multicols}{3}
	\begin{choices}
		\correctchoice \includegraphics{diversity_mcE} %EndChoice
		\choice \includegraphics{diversity_mcA} %EndChoice
		\choice \includegraphics{diversity_mcB} %EndChoice
		\choice \includegraphics{diversity_mcC} %EndChoice
		\choice \includegraphics{diversity_mcD} %EndChoice
	\end{choices}
\end{multicols}
%EndMC

%BeginMC
\question
In the experiment shown in the figure below, \textit{Balanus} was removed from the habitat shown on the left. Which of the following statements is a valid conclusion of this experiment?\\
\includegraphics[width=0.9\textwidth]{exam3_balanus}

\begin{choices}
	\choice \textit{Balanus} can survive only in the lower intertidal zone because it is unable to resist dehydration.%EndChoice
	\choice \textit{Balanus} is inferior to \textit{Chthamalus} in competing for space on rocks lower in the intertidal zone. %EndChoice
	\correctchoice The removal of \textit{Balanus} shows that the realized niche of \textit{Chthamalus} is smaller than its fundamental niche.%EndChoice
	\choice These two species of barnacle do not show competitive exclusion.%EndChoice
	\choice If \textit{Chthamalus} were removed, \textit{Balanus}'s fundamental niche would become larger.%EndChoice
\end{choices}
%EndMC


%BeginMC
\question
Which of the following is \emph{not} directly concerned with community ecology?

\begin{choices}
	\choice herbivory%EndChoice
	\choice predation%EndChoice
	\correctchoice intraspecific competition%EndChoice
	\choice resource partitioning%EndChoice
	\choice symbiosis%EndChoice
\end{choices}
%EndMC


%BeginMC
\question
The dispersion pattern for individuals in across the entire metapopulation shown in question \ref{ques:metapopulation} is most likely

\begin{choices}
\correctchoice clumped. %EndChoice
\choice uniform. %EndChoice
\choice random. %EndChoice
\choice even. %EndChoice
\choice gene flow. %EndChoice
\end{choices}
%EndMC

%BeginMC
\question
Given the following values and assuming exponential population growth, what will be the total population size after one generation has elapsed? Round to nearest whole individual.\vspace{1ex}

\hspace{\leftmargin}%
\begin{minipage}{2in}
$b = 0.09$ \qquad $N = 18,657$\\
$d = 0.12$ \qquad $K = 15,000$\\
\end{minipage}

\begin{choices}
\choice $15,000$%EndChoice
\choice $560$%EndChoice
\correctchoice $18,097$%EndChoice
\choice $19,217$ %EndChoice
\choice $3,657$ %EndChoice
\end{choices}
%EndMC

%BeginMC
\question
Which of the following is most likely to contribute to density-dependent regulation of populations?

\begin{choices}
	\choice The removal of toxic waste by decomposers.%EndChoice
	\correctchoice Intraspecific competition for nutrients.%EndChoice
	\choice Earthquakes.%EndChoice
	\choice Floods.%EndChoice
	\choice Fires.%EndChoice
\end{choices}
%EndMC

%BeginMC
\question
The ``role'' or ``job'' of a species in a community, represented by all of the resources needed by the species, is called 

\begin{choices}
	\correctchoice the niche.%EndChoice
	\choice the resource.%EndChoice
	\choice a predator.%EndChoice
	\choice a predator or prey.%EndChoice
	\choice a producer or consumer.%EndChoice
\end{choices}
%EndMC

